\section{Exact Model Properties}
\label{app:model-properties}

The source files below are authoritative for the formulas. Listings are
included directly from those files, without rewriting quantifiers or event
arguments. The two common lemmas have identical text in all three models and
are printed once; together with the eight variant-specific lemmas, they
account for all 14 model--property instances in
Table~\ref{tab:verification-results}. The model hashes in
Appendix~\ref{app:reproducibility} identify this source version.

\begin{table}[htbp]
  \centering
  \caption{Model files and property-instance counts.}
  \label{tab:appendix-model-files}
  \small
  \begin{tabularx}{\linewidth}{@{}lXrr@{}}
    \toprule
    Variant & File under \path{tamarin/rq-v2-minimal/} & Common & Specific \\
    \midrule
    Relaxed & \path{rqv2_relaxed.spthy} & 2 & 2 \\
    Message-level & \path{rqv2_message_dedup.spthy} & 2 & 3 \\
    Party-level & \path{rqv2_party_admission.spthy} & 2 & 3 \\
    \bottomrule
  \end{tabularx}
\end{table}

The source-to-abstraction mapping is given in
Table~\ref{tab:source-abstraction}. Each file has the common creation,
collection, and processing rules described in Section~4; its admission
rules select the relaxed, message, or party predicate. The guarded files
also contain the following restriction:
\begin{lstlisting}[style=modelcode]
restriction Inequality:
  "All x #i. Neq(x,x) @ #i ==> F"
\end{lstlisting}
This restricts traces containing inequality actions. It does not require a
collected pair to make progress, so it contributes no liveness guarantee.

% Keep each lemma intact across page boundaries; first blocks also keep
% their subsection heading and explanation with the first formula.
\medskip
\noindent\begin{minipage}{\linewidth}
\subsection{Common correspondence and scoped injectivity}
These all-traces formulas occur in each of the three theories. The first
requires a prior exact sender origin. The second fixes both \texttt{bid} and
\texttt{rst}; it does not forbid reuse in another batch or require every Send
to have a matching acceptance.
\lstinputlisting[style=modelcode,firstline=95,
  lastline=100]{../tamarin/rq-v2-minimal/rqv2_relaxed.spthy}
\end{minipage}\par\medskip
\noindent\begin{minipage}{\linewidth}
\lstinputlisting[style=modelcode,firstline=102,
  lastline=109]{../tamarin/rq-v2-minimal/rqv2_relaxed.spthy}
\end{minipage}\par\medskip

\noindent\begin{minipage}{\linewidth}
\subsection{Relaxed reachability}
The first existence lemma requires only one acceptance after admission. The
second requires two and a unique matching Send, without excluding unrelated
Send events.
\lstinputlisting[style=modelcode,firstline=74,
  lastline=81]{../tamarin/rq-v2-minimal/rqv2_relaxed.spthy}
\end{minipage}\par\medskip
\noindent\begin{minipage}{\linewidth}
\lstinputlisting[style=modelcode,firstline=83,
  lastline=93]{../tamarin/rq-v2-minimal/rqv2_relaxed.spthy}
\end{minipage}\par\medskip

\noindent\begin{minipage}{\linewidth}
\subsection{Message-level properties}
The rejection query does not bind its Send to the rejected batch. In contrast,
the same-party existence formula binds both accepted tuples and their two
origins to the same batch context. The final all-traces formula states
message distinction between different acceptance occurrences.
\lstinputlisting[style=modelcode,firstline=85,
  lastline=90]{../tamarin/rq-v2-minimal/rqv2_message_dedup.spthy}
\end{minipage}\par\medskip
\noindent\begin{minipage}{\linewidth}
\lstinputlisting[style=modelcode,firstline=92,
  lastline=105]{../tamarin/rq-v2-minimal/rqv2_message_dedup.spthy}
\end{minipage}\par\medskip
\noindent\begin{minipage}{\linewidth}
\lstinputlisting[style=modelcode,firstline=107,
  lastline=112]{../tamarin/rq-v2-minimal/rqv2_message_dedup.spthy}
\end{minipage}\par\medskip

\noindent\begin{minipage}{\linewidth}
\subsection{Party-level properties}
The first formula does not bind its two Sends to the rejected pair. The
second is the distinct-party non-vacuity query. The final safety formula
observes accepted occurrences; the admission-level requirement itself comes
from \texttt{AdmitDistinctParties} and the inequality restriction.
\lstinputlisting[style=modelcode,firstline=85,
  lastline=94]{../tamarin/rq-v2-minimal/rqv2_party_admission.spthy}
\end{minipage}\par\medskip
\noindent\begin{minipage}{\linewidth}
\lstinputlisting[style=modelcode,firstline=96,
  lastline=108]{../tamarin/rq-v2-minimal/rqv2_party_admission.spthy}
\end{minipage}\par\medskip
\noindent\begin{minipage}{\linewidth}
\lstinputlisting[style=modelcode,firstline=110,
  lastline=115]{../tamarin/rq-v2-minimal/rqv2_party_admission.spthy}
\end{minipage}\par\medskip
